\section{Related Literature}\label{literature}

This paper engages three distinct literatures: cartel detection via bid-pattern screens, network-based approaches to collusion identification, and the economics of inter-firm worker mobility.

\subsection{Detection through bid-pattern screens}

A foundational literature exploits the prediction that bid-rotation agreements compress the statistical properties of bids. \citet{porter1993detecting} introduced variance-based screens to the railroad freight cartel of the 1880s, showing that collusive phases exhibit lower bid dispersion than competitive phases when the cartel's internal records provide an exogenous regime switch. \citet{pesendorfer2000study} documented bid rotation in Florida school-milk procurement, demonstrating that cartel members systematically coordinate winning turns. \citet{bajari2003detecting} proposed exchangeability tests for bidder costs in sealed-bid auctions, applying them to highway paving contracts in the U.S.\ Midwest; rejections of exchangeability flag candidate collusive pairs. \citet{conley2016detecting} extended these ideas to Italian procurement, using pair-level bid-correlation tests to detect bidder groups. \citet{imhof2018screening} combined classical variance and bid-level screens with density tests in a Swiss procurement setting, demonstrating that multi-screen portfolios improve detection relative to any single indicator. \citet{abrantesmetz2006variance} provided practical guidance on implementing price-fixing screens, cautioning against single-indicator approaches. \citet{harrington2008detecting} synthesized the design principles for detection mechanisms, emphasizing that the informativeness of any screen depends on the competitive counterfactual and the cartel's coordination technology---a point our findings make concrete.

More recently, \citet{kawai2023using} proposed a regression discontinuity design centered on the winning--losing margin in Japanese construction procurement: under competition, firm characteristics should be continuous at the bid margin of zero, and any sharp discontinuity identifies a non-competitive wedge. Their approach has high internal validity but, as our results show, may lack discriminative power in settings where collusion does not concentrate at the bid margin. \citet{chassang2019collusion} offered a complementary theoretical framework in which bid constraints endogenously shape the collusive technology, predicting that auction-format-specific screens may outperform generic variance tests---a prediction consistent with our finding that classical screens behave differently across pregão and convite formats.

A common thread in this literature is that almost all validation exercises rely on ground truth from high-income economies with deep bidder pools and strong enforcement capacity. \citet{levenstein2006determines} documented the institutional determinants of cartel success and duration across 81 international cartels, noting that market structure, transparency, and enforcement intensity jointly shape collusive outcomes. What remains scarce is direct validation of off-the-shelf screens against criminal ground truth in middle-income settings with structurally different labor markets and thinner bidder pools---precisely the gap our conviction-anchored benchmark addresses.

\subsection{Network-based and machine-learning approaches}

A growing literature moves beyond bid distributions to exploit relational and network information. \citet{wachs2019network} applied co-bidding network analysis to roughly 150{,}000 Georgian procurement notices, showing that group-level cohesion and exclusivity metrics flag candidate cartels even in the absence of criminal ground truth. \citet{signor2020detection} exploited the dataset released by Brazil's \emph{Operação Lava Jato} investigation to build a probabilistic detector of collusive tenders in the Petrobras construction-cartel case, reporting accuracy between 81 and 96 percent. \citet{huber2019machine} combined classical statistical screens with supervised machine-learning classifiers and correctly predicted 84 percent of Swiss tenders out of sample, demonstrating the gains from combining bid-distribution features with flexible functional forms. \citet{clark2018bid} studied collusion and corruption jointly in Québec construction procurement, documenting how collusive networks co-evolve with political connections and how entry deterrence interacts with bid rigging.

These network approaches share a common insight: collusion is fundamentally a relational phenomenon, and screens that exploit the structure of firm--firm interactions can extract signal that univariate bid statistics miss. Our contribution builds on this insight but shifts the network channel from co-bidding patterns (observable within the procurement data itself) to worker mobility across firms (observable only in matched employer--employee administrative records). This shift is consequential because worker-flow networks are constructed from an entirely independent data source, reducing the risk of circularity between the screen and the behavioral patterns it is designed to detect.

\subsection{Worker mobility as an inter-firm transmission channel}

A separate literature in labor economics establishes that worker flows between firms serve as a vehicle for knowledge transmission, coordination, and productivity spillovers. \citet{stoyanov2012worker} showed that productivity gains propagate across Danish firms through the mobility of workers, with the magnitude of the spillover proportional to the productivity gap between sending and receiving firms. \citet{serafinelli2019good} documented similar patterns in the Veneto region of Italy, demonstrating that workers moving from high-productivity firms raise wages and productivity at destination firms. \citet{jarosch2021learning} developed an equilibrium model of learning from coworkers, showing that team composition and worker reallocation jointly determine firm-level human capital accumulation. Together, these results establish that worker flows are not merely a labor-market phenomenon but an active channel of inter-firm information transmission.

We bring this insight to the collusion-detection setting. If worker mobility transmits knowledge and facilitates coordination between firms in general, it may also serve as a coordination technology for cartel members---and, critically, leave detectable traces in administrative records. In markets with thin specialized-labor pools, where the universe of workers with the requisite expertise is small, the signal-to-noise ratio of worker-flow screens should be highest: cartel-related personnel transfers stand out against a sparse baseline of spontaneous mobility. Our specialization placebo formalizes this intuition by decomposing the worker-flow detection signal into a same-market specialization component (present even among non-colluding firms in thin markets) and a cartel-specific residual (the incremental coordination signature). To our knowledge, ours is the first paper to deploy matched employer--employee data as a cartel-detection screen and to validate it against formal antitrust convictions.
